\section{CRITICAL FLAW: ReLU Dense Span Not Proven}
\label{sec:relu-dense-span-flaw}

\subsection{The Core Issue}

\textbf{The user's insight:} If all low-rank functions converge to a function that is 0 on an interval, there's a potential problem. However, because $\text{ReLU}(af+b)$ with random $a,b$ always gives dense span, everything is fine. \textbf{But this dense span property is NOT PROVEN - it's taken for granted!}

\subsection{What's Missing}

The proof assumes that:
\[
\text{supp}(f_1(C_1)) = \mathbb{R}^d \quad \Rightarrow \quad \{\varphi_1(\langle f_1(c_1), \cdot\rangle) : c_1 \in \Omega_1\} \text{ is dense in } L^2(\mathcal{P}_X)
\]

But this implication is \textbf{never proven}. It's just stated in Assumption~\ref{assump:diversity} as if it's obvious.

\subsection{Why This Is Critical}

\begin{enumerate}
\item \textbf{The full-width case requires:}
\[
\text{supp}(w_1^0(C_1), w_2^0(C_1, \cdot)) = \mathbb{R}^d \times L^2(P_2)
\]
This is a joint support condition that ensures the function class has full expressivity.

\item \textbf{The low-rank case claims:}
\[
\text{supp}(f_1(C_1)) = \mathbb{R}^d \quad \text{(frozen random features)}
\]
And then assumes (without proof) that this gives dense span in $L^2(\mathcal{P}_X)$.

\item \textbf{The key technical crux:} For any function $f$ that is non-zero, $\text{ReLU}(af+b)$ for random $a,b$ forms a dense span. This is the crucial property that makes everything work, but it's \textbf{not proven}.
\end{enumerate}

\subsection{What Needs to Be Proven}

\textbf{Theorem (needs proof):} Let $\varphi_1$ be ReLU (or any non-polynomial activation), and let $f_1(C_1)$ be random vectors with $\text{supp}(f_1(C_1)) = \mathbb{R}^d$. Then the function class:
\[
\mathcal{F} = \{\varphi_1(\langle f_1(c_1), \cdot\rangle) : c_1 \in \Omega_1\}
\]
has dense span in $L^2(\mathcal{P}_X)$ (with respect to the data distribution $\mathcal{P}_X$).

\subsection{Why This Is Non-Trivial}

\begin{enumerate}
\item \textbf{Full support $\neq$ dense span:} Just because $\text{supp}(f_1(C_1)) = \mathbb{R}^d$ doesn't automatically mean the span is dense. You need to show that the linear combinations of these functions can approximate any function in $L^2(\mathcal{P}_X)$.

\item \textbf{ReLU-specific properties:} ReLU has special properties (piecewise linear, non-polynomial) that enable universal approximation, but this needs to be proven rigorously.

\item \textbf{Data distribution matters:} The dense span property depends on $\mathcal{P}_X$. For example, if $\mathcal{P}_X$ is supported on a lower-dimensional manifold, the dense span property might fail.

\item \textbf{Infinite-width limit:} The proof needs to work in the infinite-width limit, where $\Omega_1$ is a continuous space, not just a finite set.
\end{enumerate}

\subsection{Standard Results That Could Help}

\begin{enumerate}
\item \textbf{Leshno et al. (1993):} For non-polynomial activations, single-hidden-layer networks are dense in $C(\mathbb{R}^d)$ (and hence in $L^2$ for compactly supported measures).

\item \textbf{Cybenko (1989):} For sigmoid activations, the span is dense in $C([0,1]^d)$.

\item \textbf{Extension to ReLU:} ReLU is non-polynomial, so similar results should hold, but the proof needs to account for:
\begin{itemize}
\item The piecewise linear structure of ReLU
\item The fact that we're using random features (not trainable weights)
\item The infinite-width limit with measure-valued weights
\end{itemize}
\end{enumerate}

\subsection{The Gap in the Current Proof}

\textbf{Current state:}
\begin{itemize}
\item Assumption~\ref{assump:diversity} states: "The support of $\rho^1$ is $\mathbb{R}^d$. This ensures that the random features have dense span in $L^2(\mathcal{P}_X)$."
\item But the word "ensures" is not backed by a proof!
\item The connection between full support and dense span is assumed, not proven.
\end{itemize}

\textbf{What's needed:}
\begin{enumerate}
\item \textbf{Prove the dense span property:} Show that if $\text{supp}(f_1(C_1)) = \mathbb{R}^d$ and $\varphi_1$ is ReLU (non-polynomial), then $\mathcal{F}$ is dense in $L^2(\mathcal{P}_X)$.

\item \textbf{Handle the infinite-width case:} The proof needs to work when $\Omega_1$ is a continuous space (not just finite neurons).

\item \textbf{Account for data distribution:} Show that the dense span property holds for the specific data distribution $\mathcal{P}_X$ (not just for all distributions).

\item \textbf{Reference standard results:} Cite Leshno et al. (1993) or similar, and show how it applies to the random feature setting.
\end{enumerate}

\subsection{Why This Matters for the Proof}

If the dense span property fails, then:
\begin{itemize}
\item The universal approximation property (Assumption~\ref{assump:uap}) fails
\item The global convergence proof breaks down (it relies on being able to approximate any target function)
\item The entire theoretical framework collapses
\end{itemize}

\subsection{Solution: Adapting the Full-Width Argument}

\textbf{Key insight from the user:} We can adapt the algebraic topology argument from the full-width case (Lemma~\ref{lem:full-support-3} in \cite{nguyen2023rigorousframeworkmeanfield}, lines 2538-2672), but it's actually \textbf{simpler} for frozen features!

\textbf{The full-width case:}
\begin{itemize}
\item Uses algebraic topology to prove that full support is \emph{maintained} during training
\item Shows that the mapping $u_1 \mapsto w_1^*(t, u_1)$ is surjective using homotopy arguments
\item Then uses that $\{\varphi_1(\langle u, \cdot\rangle) : u \in \mathbb{R}^d\}$ has dense span (assumed, line 2725)
\end{itemize}

\textbf{The low-rank case (simpler!):}
\begin{itemize}
\item \textbf{No algebraic topology needed:} Since $f_1(C_1)$ are frozen, $\text{supp}(f_1(C_1)) = \mathbb{R}^d$ trivially
\item \textbf{Direct application:} We can directly prove dense span using:
\begin{enumerate}
\item Leshno et al. (1993): $\{\varphi_1(\langle u, \cdot\rangle) : u \in \mathbb{R}^d\}$ has dense span
\item Full support $\Rightarrow$ dense parameterization: $\{f_1(c_1) : c_1 \in \Omega_1\}$ is dense in $\mathbb{R}^d$
\item Continuity: $u \mapsto \varphi_1(\langle u, \cdot\rangle)$ is continuous in $L^2(\mathcal{P}_X)$
\end{enumerate}
\item \textbf{Result:} Dense parameterization + continuity $\Rightarrow$ dense span
\end{itemize}

\textbf{See `relu_dense_span_proof.tex` for the complete proof.}

\subsection{Recommendation}

\textbf{This is a CRITICAL flaw that must be fixed.} The proof needs:

\begin{enumerate}
\item A \textbf{theorem} (Theorem~\ref{thm:relu-dense-span} in `relu_dense_span_proof.tex`) proving that $\text{supp}(f_1(C_1)) = \mathbb{R}^d$ implies dense span in $L^2(\mathcal{P}_X)$ for ReLU activations.

\item A \textbf{reference} to standard universal approximation results (Leshno et al. 1993, Theorem 1).

\item A \textbf{proof} that adapts these results to:
\begin{itemize}
\item The random feature setting (frozen $f_1(C_1)$)
\item The infinite-width limit (measure-valued)
\item The specific data distribution $\mathcal{P}_X$
\end{itemize}

\item \textbf{Explicit statement} that this is a non-trivial property that requires proof, not just an assumption.

\item \textbf{Connection to full-width case:} Explain that we use a similar approach but it's simpler because features are frozen (no need for algebraic topology to maintain support).
\end{enumerate}

\subsection{Connection to User's Insight}

The user correctly identifies that:
\begin{itemize}
\item The key technical crux is showing that $\text{ReLU}(af+b)$ for random $a,b$ forms a dense span
\item This is what makes the low-rank framework work (unlike full-width which needs correlated initialization)
\item But this property is \textbf{not proven} - it's just assumed
\item This is a major gap that needs to be addressed
\end{itemize}

\textbf{This is arguably the most important theoretical gap in the entire proof.} Without proving the dense span property, the entire global convergence argument is on shaky ground.
